\chapter{Deep Dive: Service Discovery --- Eureka Internals, Failure Modes, and the Kubernetes Alternative}
\label{ch:dd-discovery}

Chapter~3 explained Eureka in four bullet points --- register, heartbeat,
cache, self-preserve --- and Chapter~27 sketched its retirement in two
lines. Between those two lies everything an operator actually meets: the
timers that decide how long a dead instance keeps receiving traffic, the
threshold that switches the registry into a mode where it stops telling
the truth, and the awkward fact that on Kubernetes the project runs two
discovery systems at once and only one of them is consulted by the
gateway. This chapter measures those things, breaks them, and then shows
the exact configuration that removes Eureka without breaking Compose.

Three questions organize the chapter:

\begin{enumerate}
  \item \textbf{What does a registry solve that DNS does not}, and why
  did the classic stack put the lookup on the \emph{client}?
  \item \textbf{How stale can the registry be} --- in seconds, with the
  arithmetic --- and what does staleness do during a deploy?
  \item \textbf{What replaces each Eureka feature} when the platform is
  Kubernetes, and how do you migrate without a flag day?
\end{enumerate}

\section{Why a registry, and why on the client}

The problem is addressing. Instances have IPs that change with every
restart; callers need a stable name. Three families of answer exist:

\begin{itemize}
  \item \textbf{Server-side load balancer.} Callers dial one fixed
  address (an nginx, an AWS ELB); it knows the instances. Simple for
  callers, but every hop adds latency and the balancer must itself be
  kept current.
  \item \textbf{Client-side registry} (Eureka, Consul). Each instance
  registers; each caller downloads the list and balances locally. No
  extra hop, no balancer to scale, callers survive a registry outage on
  their cache --- and every caller must embed a client library.
  \item \textbf{Platform DNS} (Kubernetes Services). The platform
  watches pod readiness and keeps a virtual IP's endpoint list current;
  callers use plain DNS and know nothing. No library, no registry
  process, but only inside that platform.
\end{itemize}

Netflix built the second in 2012 because it ran on EC2 with no platform
DNS to lean on and thousands of JVM callers that could carry a library.
Spring Cloud inherited it, and this project inherits it from Spring
Cloud. That history is the honest answer to ``why Eureka'': not because
it is the best fit for a Kubernetes cluster in 2026, but because the
classic stack is still what most Java shops run and what most Java
interviews ask about.

\begin{decision}[client-side, then platform-side]
The project keeps Eureka through Compose and the first Kubernetes
milestones for two reasons: Compose has no readiness-gated DNS, so a
registry is the only way for the gateway to skip a service that is up
but not ready; and running both systems side by side is how you learn
what the platform replaces. The moment Compose stops being a deployment
target --- it is already only a dev tool --- the registry is pure cost:
one more JVM, one more startup dependency, and a second opinion about
which pods are alive that disagrees with the kubelet's for up to four
minutes. Section~\ref{sec:discovery-enhance} makes the switch.
\end{decision}

\section{Reading the real configuration}

The server is a port and an annotation. All the behaviour lives in
defaults the project never sets, so read the client side first:

\begin{lstlisting}[language=yaml]
# course-service/src/main/resources/application.yml
eureka:
  client:
    service-url:
      defaultZone: http://localhost:8761/eureka/     (*@\co{1}@*)
\end{lstlisting}

\begin{lstlisting}[language=yaml]
# deploy/k8s/base/course-service.yaml (env excerpt)
- { name: EUREKA_CLIENT_SERVICEURL_DEFAULTZONE,
    value: "http://discovery-server:8761/eureka/" }  (*@\co{2}@*)
- { name: EUREKA_INSTANCE_PREFER_IP_ADDRESS, value: "true" } (*@\co{3}@*)
\end{lstlisting}

And the consumer, the gateway's routes:

\begin{lstlisting}[language=yaml]
# configurations/api-gateway.yml (excerpt)
- id: course-service-courses
  uri: lb://course-service                          (*@\co{4}@*)
  predicates:
    - Path=/api/courses/**
  filters:
    - StripPrefix=2
\end{lstlisting}

\begin{description}
  \item[\co{1}] The only address a client is told. Everything else ---
  where course-service is, where auth-service is --- comes back from
  this URL as a registry document.
  \item[\co{2}] On Kubernetes \code{discovery-server} is resolved by
  \emph{platform DNS} to a ClusterIP. The registry itself is found by
  the mechanism it is about to duplicate.
  \item[\co{3}] Register the pod IP, not the pod name (Chapter~3). The
  registered address is what the gateway will dial, bypassing the
  \code{course-service} Service entirely.
  \item[\co{4}] \code{lb://} hands the name to Spring Cloud LoadBalancer,
  which asks the Eureka client's local cache for instances and picks one
  round-robin. The Kubernetes Service \code{course-service:8082} is never
  consulted on this path.
\end{description}

Point \co{4} is the sentence to remember from this chapter. On the
cluster, \emph{readiness probes protect the Kubernetes Service, and the
gateway does not use the Kubernetes Service}. The only thing between
the gateway and a pod that is not ready is the Eureka status --- which
the pod reports about itself.

\subsection{The timers}

Everything below is a default the project inherits. Each is a property
you can set; each has a reason for its value.

\begin{center}
\begin{tabular}{@{}p{0.44\linewidth}rp{0.36\linewidth}@{}}
\toprule
\textbf{Property (\code{eureka.*})} & \textbf{Default} & \textbf{Meaning} \\
\midrule
\code{instance.lease-renewal-interval-in-seconds} & 30\,s & heartbeat period \\
\code{instance.lease-expiration-duration-in-seconds} & 90\,s & no heartbeat for this long = expired \\
\code{server.eviction-interval-timer-in-ms} & 60\,s & how often the server sweeps expired leases \\
\code{server.response-cache-update-interval-ms} & 30\,s & server's read-only registry cache refresh \\
\code{client.registry-fetch-interval-seconds} & 30\,s & client re-downloads the registry \\
\code{spring.cloud.loadbalancer.cache.ttl} & 35\,s & LoadBalancer's own instance-list cache \\
\code{server.renewal-percent-threshold} & 0.85 & self-preservation trigger \\
\bottomrule
\end{tabular}
\end{center}

Now the arithmetic an interviewer wants.

\textbf{A new instance receiving traffic.} Spring Cloud registers as
soon as the web server starts (the raw Netflix client would wait 40\,s;
Spring's \code{EurekaAutoServiceRegistration} forces an on-demand
registration). The server's read-only cache may be up to 30\,s stale, the
gateway's registry fetch up to 30\,s behind that, and its LoadBalancer
cache up to 35\,s behind that. Worst case: about
\textbf{95 seconds} between ``Tomcat is listening'' and ``the gateway
knows''. Typical: half that.

\textbf{A dead instance no longer receiving traffic.} The pod is killed
without a graceful shutdown (OOM, node loss). The lease expires after
90\,s without a heartbeat; the eviction sweep runs every 60\,s, so up to
150\,s after death the server still lists it; add the same 30 + 30 +
35\,s of caches: up to \textbf{245 seconds} --- four minutes --- during
which the gateway keeps dialing an IP that no longer answers. A graceful
shutdown sends a DELETE on the way out and skips the first 150\,s, leaving
up to 95\,s of stale caches.

\textbf{Self-preservation.} Every minute the server compares heartbeats
received against heartbeats expected: instances $\times$ 2 (one per
30\,s) $\times$ 0.85. Below the line, it stops evicting anything and the
dashboard shows the red banner. The project registers five clients:
expected $5 \times 2 = 10$, threshold $\lfloor 10 \times 0.85 \rfloor =
8$, and the condition is ``more than 8''. Kill \emph{one} service and
the count drops to 8 --- not more than 8 --- and the registry enters
self-preservation. On this stack, a single crashed service means
\emph{no} instance is ever evicted again until it comes back. Netflix
tuned that threshold for hundreds of instances, where losing 15\,\% in a
minute really does mean a network partition; on five, it means Tuesday.

\begin{handson}[watch a dead instance stay alive]
In Compose, stop course-service without letting it deregister, then
ask the registry about it a minute later:
\begin{lstlisting}[language=bash]
$ docker compose kill course-service
$ sleep 60
$ curl -s -H 'Accept: application/json' \
    http://localhost:8761/eureka/apps/COURSE-SERVICE | \
    python -c "import sys,json; a=json.load(sys.stdin)['application']; \
      print([i['status'] for i in a['instance']])"
['UP']
\end{lstlisting}
Now open \code{http://localhost:8761}. After the next sweep, the
instance is either gone --- or, because four remaining clients send 8
renewals a minute, the page carries the banner \textsc{emergency! eureka
may be incorrectly claiming instances are up when they're not}, and
course-service stays \code{UP} for as long as you care to wait. Meanwhile
\code{curl localhost:8080/api/courses} returns a 503 from the gateway on
every attempt that picks the dead address.
\end{handson}

\section{Pros and cons, honestly}

\begin{itemize}
  \item \textbf{For.} No extra network hop; callers keep working through
  a registry outage on their cache; one mechanism that is identical on a
  laptop, in Compose, and in the cluster; a dashboard that shows who is
  registered; and an instance can mark itself \code{OUT\_OF\_SERVICE}
  for a controlled drain.
  \item \textbf{Against, in this project.} A second source of truth about
  liveness that lags the kubelet's by up to four minutes; readiness
  probes bypassed on the gateway path (\co{4} above); self-preservation
  tuned for a fleet a hundred times larger; a 512\,Mi JVM that does
  nothing for traffic; \code{prefer-ip-address} as a platform-specific
  workaround every manifest must carry; and every service compiled
  against a client library whose upstream is in maintenance mode.
  \item \textbf{Neutral.} Eureka is AP by design: it prefers answering
  with stale data to not answering. That is the right choice for a
  registry --- a caller with a slightly wrong list can retry; a caller
  with no list can do nothing --- and it is also why every number above
  is ``up to''.
\end{itemize}

\section{What breaks}

\begin{pitfall}[rolling deploy, thirty seconds of 503s]
Symptom: \code{kubectl rollout restart deployment/course-service}
succeeds, readiness passes, and for half a minute the gateway returns
\code{503 Service Unavailable} on \code{/api/courses}. Cause: the old
pod deregistered on SIGTERM and exited; the new pod registered with a
new IP; but the gateway's LoadBalancer still holds the old IP for up to
95\,s of cache. The Kubernetes Service switched instantly and nobody
asked it. Fix, in ascending cost: a \code{preStop} hook that sleeps
20--30\,s so the old pod keeps answering while caches drain
(\code{terminationGracePeriodSeconds} must exceed it); shorter caches
(\code{registry-fetch-interval-seconds: 5},
\code{spring.cloud.loadbalancer.cache.ttl: 5s}) at the price of
registry load; or a gateway \code{Retry} filter on connection errors so
the second attempt lands on the live IP. Or remove the registry from the
path (Section~\ref{sec:discovery-enhance}) and let the Service do what it
is for.
\end{pitfall}

\begin{pitfall}[registered, UP, and returning errors]
Symptom: the dashboard shows course-service \code{UP}; requests fail
with database errors. Cause: Eureka status is self-reported, and by
default it is \code{UP} the moment registration succeeds --- it does not
consult the actuator health. Flyway may have failed after registration,
Postgres may have gone away, and the instance still says \code{UP}
every 30\,s. Fix: \code{eureka.client.healthcheck.enabled: true}
propagates actuator health to the registry status, so a \code{DOWN}
health check becomes \code{DOWN} in Eureka on the next heartbeat --- 30\,s
later, plus caches. Kubernetes readiness would have pulled the pod in
5\,s; on the gateway path, it did not get the chance.
\end{pitfall}

\begin{pitfall}[two Eureka replicas, two answers]
Symptom: after scaling discovery-server to two replicas for
availability, some callers see an instance the others do not. Cause:
peers replicate registrations to each other asynchronously; a client
whose \code{defaultZone} lists both addresses talks to the first that
answers and may read a replica that has not yet received the
replication. Each server also counts renewals independently, so one may
enter self-preservation while the other evicts. Fix: list both peers in
every client's \code{defaultZone} and in each server's
\code{defaultZone} (they must register with \emph{each other}), accept
seconds of inconsistency as the price of AP, and never scale a registry
you are about to remove.
\end{pitfall}

\section{What breaks at 3\,a.m.}

The discovery-server pod dies at 3\,a.m.; nothing else does.

\begin{itemize}
  \item \textbf{Gateway traffic: unaffected.} The gateway's Eureka client
  keeps its last registry in memory and the LoadBalancer keeps routing
  to it. The registry fetch fails every 30\,s and logs a warning; the
  cache does not expire. This is the AP promise kept.
  \item \textbf{Heartbeats: fail silently.} Every client logs a renewal
  failure each 30\,s. Nothing else changes; leases are a server concept
  and the server is gone.
  \item \textbf{A service restarted during the outage: invisible.} It
  boots, registration fails, it retries in the background --- and it
  serves nothing through the gateway until the registry returns and the
  95\,s of caches drain. Its Kubernetes Service has it as an endpoint
  within seconds; nobody uses that Service.
  \item \textbf{A service that died during the outage: still routed to.}
  No server, no eviction. The gateway keeps the dead IP until the
  registry is back \emph{and} has evicted it \emph{and} the caches have
  turned over: the full 245\,s, starting from when the server returns.
  \item \textbf{When discovery-server returns: a cold registry.} It
  starts empty. For its first \code{wait-time-in-ms-when-sync-empty}
  (5 minutes by default) it refuses to serve a registry it considers
  incomplete, returning what it has; clients re-register on their next
  heartbeat within 30\,s. Full recovery: about a minute after the pod is
  Ready.
\end{itemize}

Compare with the platform path: a Service's endpoint list is maintained
by the control plane, which is replicated and which every other
component already depends on. There is no separate 3\,a.m. story for
Kubernetes DNS because it has no separate process to lose.

\section{Enhancing the resolution}
\label{sec:discovery-enhance}

The goal: on Kubernetes, route by Service DNS and let readiness probes
gate traffic; in Compose, keep Eureka so the developer experience is
unchanged. A Spring profile splits the two, and the config-server serves
the profile document on top of the base one.

\subsection{Step 1: a profile document for the gateway}

The base \code{api-gateway.yml} keeps its \code{lb://} routes. A new
\code{api-gateway-k8s.yml} beside it replaces them --- the whole list,
because YAML lists override by index and a partial override is a trap:

\begin{lstlisting}[language=yaml]
# configurations/api-gateway-k8s.yml -- served when profile k8s is active
spring:
  cloud:
    gateway:
      routes:
        - id: auth-service
          uri: http://auth-service:8081               (*@\co{1}@*)
          predicates:
            - Path=/api/auth/**
          filters:
            - StripPrefix=2
        - id: auth-oauth2
          uri: http://auth-service:8081
          predicates:
            - Path=/oauth2/**
        - id: auth-oauth2-callback
          uri: http://auth-service:8081
          predicates:
            - Path=/login/oauth2/**
        - id: course-service
          uri: http://course-service:8082
          predicates:
            - Path=/api/courses/**,/api/students/**,/api/assignments/**,/api/enrollments/**,/api/submissions/**
          filters:
            - StripPrefix=2
        - id: dictionary-service
          uri: http://dictionary-service:8083
          predicates:
            - Path=/api/dictionary/**
          filters:
            - StripPrefix=2
      httpclient:
        connect-timeout: 2000                          (*@\co{2}@*)
        response-timeout: 10s
eureka:
  client:
    enabled: false                                     (*@\co{3}@*)
\end{lstlisting}

\begin{description}
  \item[\co{1}] A Service name and port. CoreDNS resolves it to the
  ClusterIP; kube-proxy balances across Ready endpoints. Readiness
  probes are now on the gateway's path.
  \item[\co{2}] With no registry to consult, the gateway's only defence
  against a dead endpoint is a fast connect timeout; kube-proxy will
  have removed the endpoint within seconds anyway.
  \item[\co{3}] The Eureka client stays on the classpath but does
  nothing: no registration, no fetch, no heartbeat warnings. This is
  what lets one image serve both environments.
\end{description}

\subsection{Step 2: Feign without the registry}

course-service's \code{AuthServiceClient} resolves \code{auth-service}
through the LoadBalancer. Give it an optional fixed URL:

\begin{lstlisting}[language=Java]
// course-service/.../client/AuthServiceClient.java
@FeignClient(
        name = "auth-service",
        url = "${app.clients.auth-service:}",   // empty -> use lb://
        fallback = AuthServiceFallback.class)
public interface AuthServiceClient {
    @GetMapping("/users/{id}")
    UserDto getUserById(@PathVariable Long id);
}
\end{lstlisting}

When \code{url} resolves to an empty string Feign falls back to
load-balanced resolution by \code{name}; when it resolves to
\code{http://auth-service:8081} the LoadBalancer is bypassed. The
profile document for course-service sets it:

\begin{lstlisting}[language=yaml]
# configurations/course-service-k8s.yml
app:
  clients:
    auth-service: http://auth-service:8081
eureka:
  client:
    enabled: false
\end{lstlisting}

\subsection{Step 3: the manifests}

The diff against \code{deploy/k8s/base/api-gateway.yaml}; every service
manifest changes the same way:

\begin{lstlisting}
 env:
+  - { name: SPRING_PROFILES_ACTIVE, value: "k8s" }
   - { name: SPRING_CONFIG_IMPORT,
       value: "configserver:http://config-server:8888" }
-  - { name: EUREKA_CLIENT_SERVICEURL_DEFAULTZONE,
-      value: "http://discovery-server:8761/eureka/" }
-  - { name: EUREKA_INSTANCE_PREFER_IP_ADDRESS, value: "true" }
\end{lstlisting}

And in \code{kustomization.yaml}:

\begin{lstlisting}
 resources:
   - config-server.yaml
-  - discovery-server.yaml
   - api-gateway.yaml
\end{lstlisting}

Compose is untouched: no profile, the base documents, \code{lb://},
Eureka. The same images run in both. When Compose is finally retired,
delete \code{spring-cloud-starter-netflix-eureka-client} from every
\code{pom.xml}, the \code{discovery-server} module, and the base
documents' \code{lb://} routes --- in that order, each a separate commit
the pipeline can verify.

\begin{handson}[prove the platform path]
After applying the profile, confirm the gateway no longer dials pod IPs
directly --- it dials the Service, whose endpoints follow readiness:
\begin{lstlisting}[language=bash]
$ kubectl -n ts get endpoints course-service
NAME             ENDPOINTS         AGE
course-service   10.42.0.57:8082   3d
$ kubectl -n ts rollout restart deployment/course-service
$ while true; do curl -s -o /dev/null -w '%{http_code}\n' \
    -H "Authorization: Bearer $TOKEN" \
    http://ts.local/api/courses; sleep 1; done
200
200
200   # no 503 window: the endpoint switched with readiness
\end{lstlisting}
Run the same loop before the change and count the 503s.
\end{handson}

\section{Outside the project}

Consul is the registry you meet outside the JVM: an agent per host runs
health checks itself instead of trusting self-reported heartbeats, and
serves the result over DNS as well as HTTP, so callers need no library.
Service meshes (Istio, Linkerd) go one step further and put a sidecar
proxy next to every pod that does discovery, retries, timeouts, and
mTLS, leaving the application with \code{http://course-service}. In
Spring, \code{spring-cloud-kubernetes} lets \code{lb://} names resolve
against the Kubernetes API instead of Eureka --- useful during a
migration, unnecessary at its end. And AWS Cloud Map, the managed
descendant of Eureka's idea, exists mostly for ECS, where there is no
kube-proxy to make it redundant.

\section{Questions from real interviews}

\subsection*{Q: Your gateway keeps sending traffic to an instance that
died ninety seconds ago. Why, and how do you cut that down?}

Because nothing on the path noticed. Eureka learns of an ungraceful
death only by a missing heartbeat: 90\,s lease expiry, then a sweep that
runs every 60\,s, then a 30\,s server cache, a 30\,s client fetch, and a
35\,s LoadBalancer cache --- up to four minutes end to end, and if the
registry has slipped into self-preservation, forever. First cut the
tail: shorten \code{registry-fetch-interval-seconds} and the
LoadBalancer TTL to 5\,s, which costs registry traffic you can afford
below a few hundred instances. Then cut the head: make death graceful
(SIGTERM handling, \code{terminationGracePeriodSeconds} long enough) so
the DELETE goes out and the 150\,s of lease logic never runs. Then make
the caller resilient regardless: a \code{Retry} filter on connect
failures, so a stale address costs one retry, not an error. The senior
answer ends with ``or stop routing by registry on Kubernetes, because
kube-proxy removes the endpoint within seconds of the readiness probe
failing''.

\subsection*{Q: How does a new instance start receiving traffic during a
rolling deploy with zero errors?}

Two conditions and an overlap. The new instance must not receive
traffic before it can serve: on Kubernetes that is the readiness
probe, on Eureka it is registering only after the context is fully up
--- Spring does this, but a health check that passes before Flyway
finishes would not, so probe something that depends on the database.
The old instance must keep serving until every caller has stopped
sending: on Kubernetes the endpoint is removed at SIGTERM but in-flight
requests finish inside the grace period, and a \code{preStop} sleep
covers callers with stale caches. The overlap is
\code{maxSurge: 1, maxUnavailable: 0}: bring the new one up, wait for
Ready, then take the old one down. Errors during a rollout are almost
always one of three things --- a probe that lies, a cache that
outlives the pod, or a grace period shorter than the slowest request
--- and you name which one from the shape of the error.

\subsection*{Q: Eureka versus Kubernetes DNS. One minute.}

Both map a name to a set of addresses; they differ in who maintains the
set and who does the lookup. Eureka: the instances register themselves,
report their own health, and every caller carries a client that
downloads and caches the list --- no platform needed, works on VMs,
AP with minutes of staleness, and every caller is a JVM with a
library. Kubernetes: the control plane maintains a Service's endpoints
from label selectors and readiness probes it runs itself, kube-proxy
balances, callers use DNS and know nothing --- seconds of staleness,
no library, no registry process, but only inside the cluster. On
Kubernetes, Eureka is a second, slower opinion about liveness. Keep it
only if you must run the same code somewhere the platform is not.

\subsection*{Q: You run two Eureka replicas and they disagree about
whether an instance is alive. What now?}

Expect it; design for it. Peer replication is asynchronous and each
replica evicts and self-preserves on its own counters, so a window of
disagreement is inherent --- Eureka is AP and the disagreement is the
cost of answering during a partition. What matters is that clients
tolerate it: every client lists both replicas, retries the second on
failure, and treats the registry as a hint that a connect timeout can
override. If the disagreement persists past a few minutes, one replica
is in self-preservation and the other is not; check the renewal
threshold against the actual instance count, and on a small fleet
lower \code{renewal-percent-threshold} or disable self-preservation
outright. The wrong answer is trying to make the registry consistent;
the right one is making callers indifferent to which replica they read.

\subsection*{Q: Two hundred services, each fetching the registry every
thirty seconds. The registry is melting. What do you change?}

Four hundred fetches a minute of a document listing two hundred
applications is the first bottleneck, and Eureka's answer is delta
fetches: clients ask for \code{/apps/delta} after the first full
download, which is on by default and only fails when the server's
delta queue is disabled. Second, the read-only response cache exists
precisely so the server serializes the registry once per 30\,s, not
once per request --- keep it. Third, split by zone: Eureka's
\code{availability-zone} configuration lets clients prefer a local
registry and a local instance, so each server carries a fraction of
the fetches. Beyond a thousand instances, Netflix's own answer was
more Eureka servers behind DNS, and today's answer is that at that
size you are on a platform that does discovery for you.

\begin{quiz}
\begin{enumerate}
  \item Compute the worst-case delay before the gateway routes to a
  freshly started course-service pod, naming each cache in the chain
  and its default.
  \item An instance is killed by the OOM killer. List the timers that
  must elapse before the gateway stops dialing it, and state which of
  them a graceful shutdown would skip.
  \item With five registered clients, show why the loss of one puts the
  registry into self-preservation, and give the one-line configuration
  that changes that on a small fleet.
  \item On Kubernetes, which mechanism gates traffic to a pod on the
  gateway path today, and why do readiness probes not? What single
  configuration change puts readiness back on the path?
  \item After \code{eureka.client.enabled: false} under the \code{k8s}
  profile, explain how the same course-service image still registers
  with Eureka in Compose, and what \code{url = "\${app.clients.auth-service:}"}
  does in each environment.
  \item Interview: ``Why does a registry keep serving a dead instance
  while a Kubernetes Service does not?'' Answer in terms of who
  observes liveness and how.
\end{enumerate}
\end{quiz}
